\documentclass[11pt,a4paper]{article}

% ─────────────────────────────────────────────────────────────────────────────
%  Cobertura de Rúbricas (Parcial + Final) — Sistema Steam Distribuido
%  Documento de trazabilidad: qué se hizo, por qué y dónde se ve en el código.
%  Compilar con: pdflatex Cobertura_Rubricas_Parcial_Final.tex   (x2 para el índice)
% ─────────────────────────────────────────────────────────────────────────────

\usepackage[utf8]{inputenc}
\usepackage[T1]{fontenc}
\usepackage[spanish,es-noshorthands]{babel}
\usepackage[margin=2.2cm]{geometry}
\usepackage{xcolor}
\usepackage{enumitem}
\usepackage{longtable}
\usepackage{listings}
\usepackage{titlesec}
\usepackage{fancyhdr}
\usepackage[colorlinks=true,linkcolor=azul,urlcolor=azul]{hyperref}

% ── Colores ──────────────────────────────────────────────────────────────────
\definecolor{azul}{RGB}{20,80,140}
\definecolor{acento}{RGB}{180,60,20}
\definecolor{verde}{RGB}{30,120,60}
\definecolor{gris}{RGB}{120,120,120}
\definecolor{grisclaro}{RGB}{245,245,245}

% ── Estilo de código ─────────────────────────────────────────────────────────
\lstdefinestyle{java}{
  basicstyle=\ttfamily\footnotesize,
  breaklines=true,
  frame=single,
  rulecolor=\color{gris},
  backgroundcolor=\color{grisclaro},
  keywordstyle=\color{azul}\bfseries,
  commentstyle=\color{verde}\itshape,
  stringstyle=\color{acento},
  language=Java,
  showstringspaces=false,
  numbers=left,
  numberstyle=\tiny\color{gris},
  tabsize=2,
  literate={í}{{\'i}}1 {á}{{\'a}}1 {é}{{\'e}}1 {ó}{{\'o}}1 {ú}{{\'u}}1 {ñ}{{\~n}}1
}
\lstset{style=java}

% ── Encabezado / pie ─────────────────────────────────────────────────────────
\pagestyle{fancy}
\fancyhf{}
\lhead{\small Sistema Steam Distribuido}
\rhead{\small Cobertura de Rúbricas}
\cfoot{\thepage}
\renewcommand{\headrulewidth}{0.4pt}

% ── Macros para cada punto de rúbrica ────────────────────────────────────────
\newcommand{\Que}{\par\smallskip\noindent\textbf{\textcolor{azul}{¿Qué se hizo?}}\quad}
\newcommand{\Por}{\par\smallskip\noindent\textbf{\textcolor{azul}{¿Por qué?}}\quad}
\newcommand{\Don}{\par\smallskip\noindent\textbf{\textcolor{azul}{¿Dónde se ve?}}\quad}
\newcommand{\nivel}[1]{\hfill\fcolorbox{verde}{white}{\small\textcolor{verde}{\textbf{#1}}}}
\newcommand{\pend}[1]{\hfill\fcolorbox{acento}{white}{\small\textcolor{acento}{\textbf{#1}}}}
\newcommand{\codigo}[1]{\texttt{#1}}

\titleformat{\subsubsection}{\normalfont\large\bfseries\color{azul}}{\thesubsubsection}{1em}{}

\title{\textbf{Sistema Steam Distribuido}\\[4pt]
\large Cobertura punto por punto de las rúbricas\\ Parcial y Final\\[6pt]
\normalsize ICI-4344 — Computación Paralela y Distribuida}
\author{Documento de trazabilidad Rúbrica $\rightarrow$ Código}
\date{\today}

\begin{document}
\maketitle
\thispagestyle{fancy}

\begin{center}
\fcolorbox{azul}{grisclaro}{\parbox{0.92\textwidth}{\small
\textbf{Cómo leer este documento.} Por cada criterio de las dos rúbricas se responde:
\textcolor{azul}{\textbf{¿Qué se hizo?}} (la implementación concreta),
\textcolor{azul}{\textbf{¿Por qué?}} (la justificación técnica) y
\textcolor{azul}{\textbf{¿Dónde se ve?}} (archivo / clase / método donde está cubierto).
La etiqueta \textcolor{verde}{\textbf{verde}} marca lo ya implementado y verificado en código;
la \textcolor{acento}{\textbf{naranja}} marca lo que queda como trabajo de informe o presentación.}}
\end{center}

\tableofcontents
\newpage

% ═════════════════════════════════════════════════════════════════════════════
\section{Resumen del sistema}
% ═════════════════════════════════════════════════════════════════════════════

\textbf{Dominio.} Plataforma distribuida tipo \emph{Steam}: catálogo de videojuegos con
compra-venta y mensajería entre usuarios. Es la evolución del proyecto parcial
(cliente--servidor) a una arquitectura multinodo con coordinación distribuida.

\textbf{Dos funciones principales de gran escala.}
\begin{enumerate}[leftmargin=1.4em]
  \item \textbf{F1 — Compra de juegos (transacción en 2 fases):} reserva temporal con
        TTL de 5 min sobre \emph{stock finito}, y confirmación de pago contra billetera.
        Es el recurso crítico protegido por exclusión mutua.
  \item \textbf{F2 — Mensajería 1-a-1:} envío con almacenamiento \emph{offline} y
        ordenamiento causal de la conversación mediante relojes de Lamport.
\end{enumerate}

\textbf{Topología (7 procesos / JVM independientes).}

\begin{center}
\small
\begin{tabular}{l l l}
\hline
\textbf{Componente} & \textbf{Rol} & \textbf{Puerto(s)} \\
\hline
Proxy            & Gateway: ruteo Round-Robin, health-check, membresía & 8080 \\
svSesiones 1 / 2 & Autenticación y tokens de sesión & 8081 / 8181 \\
svJuegos 1 / 2   & Transacciones (F1) + Bully + Mutex & 8082 / 8282 \\
svMensajeria 1/2 & Chat (F2) con orden causal & 8083 / 8383 \\
(canal Bully)    & Elección de coordinador entre nodos de juegos & 9082 / 9282 \\
(canal Mutex)    & Exclusión mutua del recurso \codigo{stock} & 9182 / 9382 \\
\hline
\end{tabular}
\end{center}

\textbf{Stack.} Java (compila con \codigo{javac --release 17}), sockets TCP, serialización
JSON con Gson 2.10.1. Sin Maven: el script \codigo{scripts/1\_build.ps1} descarga Gson,
compila los 35 archivos y genera \codigo{sistema-steam.jar}.

\textbf{Estructura de paquetes} (\codigo{src/main/java/com/steam/}):
\begin{itemize}[leftmargin=1.4em,nosep]
  \item \codigo{common/} — \codigo{Constantes}, \codigo{MensajeProtocolo} (sobre JSON),
        \codigo{GestorPersistencia}, \codigo{ValidadorToken}, \codigo{Utils},
        \codigo{RelojLamport}, \codigo{GestorLog}, \codigo{GestorSnapshot}, \codigo{RegistradorProxy}.
  \item \codigo{models/} — \codigo{Usuario}, \codigo{Sesion}, \codigo{Juego}, \codigo{Reserva},
        \codigo{Venta}, \codigo{Mensaje}, y las BD en memoria \codigo{BDSesiones/BDJuegos/BDMensajeria}.
  \item \codigo{proxy/} — \codigo{Proxy}.
  \item \codigo{servidores/} — \codigo{svSesiones}, \codigo{svJuegos}, \codigo{svMensajeria}, \codigo{GestorLocks}.
  \item \codigo{coordinacion/} — \codigo{GestorBully}, \codigo{GestorMutexCentralizado},
        \codigo{RegistroMembresia}, \codigo{NodoInfo}, \codigo{EstadoNodo}.
  \item \codigo{carga/} — \codigo{GeneradorCarga}, \codigo{FallaInducida}.
  \item \codigo{watchdog/} — \codigo{WatchdogServidor}.
\end{itemize}

% ═════════════════════════════════════════════════════════════════════════════
\section{Rúbrica PARCIAL — cobertura punto por punto}
% ═════════════════════════════════════════════════════════════════════════════

\subsection{I. Informe Técnico e Inspección de Modelos (40\%)}

\subsubsection{1.1 Fundamentación y Teoría — concurrencia, fallos, transparencia}
\nivel{Base técnica en código}

\Que Se aplican los tres conceptos clave al sistema. \emph{Concurrencia:} cada servidor
atiende muchos clientes a la vez con un pool de hilos y protege los datos compartidos.
\emph{Fallos independientes:} cada nodo puede caer por separado sin tumbar al resto.
\emph{Transparencia:} el cliente solo conoce el Proxy (puerto 8080) y no sabe qué nodo
físico lo atiende.

\Por Son las tres propiedades intrínsecas de un sistema distribuido (concurrencia,
ausencia de reloj global y fallos parciales). El informe debe \emph{vincularlas al código},
no solo definirlas.

\Don Concurrencia: \codigo{Executors.newFixedThreadPool} en \codigo{svJuegos.escuchar()}.
Transparencia de ubicación: \codigo{Proxy.rutear()} (el cliente nunca pone IP de un nodo).
Transparencia de acceso: protocolo uniforme \codigo{MensajeProtocolo} para toda operación.
\pend{Diagramas/redacción: informe}

\subsubsection{1.2 Modelado de Ingeniería — físico, arquitectónico y de funciones}
\nivel{Descrito; faltan diagramas formales}

\Que \emph{Modelo físico:} 7 procesos sobre red local (LAN), comunicados por TCP.
\emph{Modelo arquitectónico:} cliente $\rightarrow$ Proxy (gateway) $\rightarrow$ clústeres de
servidores por servicio (3 capas: cliente, intermediación, servicios + persistencia).
\emph{Modelo de funciones:} F1 (compra en 2 fases) y F2 (mensajería) descritas con su
paso de mensajes.

\Por El modelado comunica el diseño antes que el código y permite razonar la arquitectura
(es más barato corregir un diagrama que reescribir código).

\Don La materia prima de los tres modelos está en este documento (Sección 1 y tablas de
puertos) y en los encabezados Javadoc de \codigo{Proxy.java} y los \codigo{sv*.java}.
\pend{Diagramas UML de secuencia: informe}

\subsubsection{1.3 Análisis Fundamental — modelo de seguridad y modelo de fallos}
\nivel{Implementado en código}

\Que \emph{Seguridad:} contraseñas con hash SHA-256, tokens de sesión UUID, validación de
token en toda operación sensible y control de acceso por rol (COMPRADOR/VENDEDOR/ADMIN).
\emph{Fallos:} clasificación Crash (nodo caído) y Omisión (mensaje perdido) con detección
por timeouts/health-check y recuperación por failover y respaldo.

\Por El canal TCP es inseguro (interceptación, suplantación) y los fallos son parciales:
hay que autenticar/autorizar y tener estrategia explícita por tipo de fallo.

\Don Seguridad: \codigo{Utils.hashPassword()} (SHA-256), \codigo{svSesiones.login()/validarToken()},
\codigo{ValidadorToken}, chequeos de rol en \codigo{svJuegos.publicarJuego()/agregarSaldo()}.
Fallos: \codigo{Proxy.iniciarHealthCheck()} (Crash) y reintento al espejo en
\codigo{Proxy.rutear()} (Omisión).

\subsection{II. Implementación y Código Fuente en Java (30\%)}

\subsubsection{2.1 Comunicación y Marshalling (10\%)}
\nivel{Excelente}

\Que Toda la comunicación es por \emph{sockets TCP}. Las estructuras complejas
(objetos, listas, mapas anidados) se serializan a JSON con Gson antes de viajar.

\Por La rúbrica pide sockets + \emph{marshalling de estructuras complejas}, no solo texto
plano. El JSON resuelve diferencias de representación entre procesos y permite enviar
estados ricos (catálogo, reservas, billeteras) en un solo mensaje.

\Don \codigo{MensajeProtocolo.toJson()/fromJson()} (sobre Gson) y su \codigo{Map<String,Object> payload}.
Sockets: \codigo{ServerSocket}/\codigo{Socket} en \codigo{Proxy.escuchar()},
\codigo{Proxy.reenviar()} y \codigo{sv*.escuchar()}.

\begin{lstlisting}[caption={Sobre de comunicación serializado a JSON (com.steam.common.MensajeProtocolo).}]
private static final Gson GSON = new Gson();
public String toJson()                       { return GSON.toJson(this); }
public static MensajeProtocolo fromJson(String j){ return GSON.fromJson(j, MensajeProtocolo.class); }
\end{lstlisting}

\subsubsection{2.2 Gestión de Concurrencia (10\%)}
\nivel{Excelente}

\Que Servidor multihilo con \emph{pool de hilos} (30 por componente) y sincronización en
las regiones críticas: \codigo{synchronized} sobre la BD, \codigo{Semaphore} para el recurso
\codigo{stock}, y estructuras \emph{lock-free} (\codigo{AtomicLong}, \codigo{AtomicInteger},
\codigo{ConcurrentHashMap}, \codigo{CopyOnWriteArrayList}).

\Por Un servidor iterativo se bloquea con un cliente lento. El pool da escalabilidad, y la
sincronización evita condiciones de carrera (p.ej. que dos compradores reserven el último
ejemplar del stock).

\Don Pool: \codigo{Executors.newFixedThreadPool(Constantes.POOL\_SIZE)} en cada servidor.
Región crítica: bloque \codigo{synchronized(lock)} en \codigo{svJuegos.comprarJuego()}.
Round-Robin atómico: \codigo{AtomicInteger} en \codigo{Proxy.rutear()}.

\subsubsection{2.3 Lógica de las Funciones Principales (10\%)}
\nivel{Excelente}

\Que Las dos funciones operan correctamente y el sistema resiste excepciones de red.
F1: \codigo{COMPRAR\_JUEGO} (reserva + descuento de stock) y \codigo{CONFIRMAR\_PAGO}
(cobro a billetera, registro de venta). F2: \codigo{ENVIAR\_MENSAJE}, \codigo{VER\_MENSAJES}
(entrega offline), \codigo{VER\_CONVERSACION}.

\Por La rúbrica exige dos funciones de gran escala \emph{y} manejo de excepciones para que
un fallo parcial de red no tumbe todo el sistema.

\Don F1: \codigo{svJuegos.comprarJuego()} y \codigo{confirmarPago()}; expiración de reservas
en \codigo{GestorLocks}. F2: \codigo{svMensajeria.enviarMensaje()/verMensajes()}.
Excepciones: \codigo{try/catch} de \codigo{IOException} y \codigo{SocketTimeoutException}
en todos los \codigo{manejarCliente()}, con failover en \codigo{Proxy.rutear()}.

\subsection{III. Presentación, Demo y Defensa (30\%)}

\subsubsection{3.1 Exposición y Video — 3.2 Defensa Técnica}
\pend{Pendiente: presentación}

\Que La \emph{capacidad} de demo está lista: se puede levantar el sistema completo y
ejercitar ambas funciones en vivo. La coherencia informe$\leftrightarrow$código está
respaldada por este documento.

\Por La nota de presentación es obligatoria e independiente de la entrega.

\Don Scripts de ejecución \codigo{scripts/2..9}. Material de defensa: este documento y
los Javadoc de cada clase.
\pend{Grabar video $\le$3 min + defensa oral: equipo}

% ═════════════════════════════════════════════════════════════════════════════
\section{Rúbrica FINAL — cobertura punto por punto}
% ═════════════════════════════════════════════════════════════════════════════

\subsection{§2 Requisitos Técnicos Distribuidos Obligatorios}

\subsubsection{2.1 Topología multinodo}
\nivel{Excelente}

\Que Tres o más nodos en procesos/JVM separados (7 en total), arquitectura
\emph{multiservidor} (clústeres por servicio) donde la coordinación NO está centralizada en
un único servidor, y \emph{descubrimiento de membresía}: cada servidor se registra solo en
el Proxy al arrancar y se da de baja al apagarse; el estado vivo se persiste.

\Por El final exige pasar de cliente--servidor a multinodo: ningún nodo solo debe dejar el
sistema fuera de servicio, y los nodos deben conocerse por una lista de membresía.

\Don Registro dinámico: \codigo{RegistradorProxy.registrarAsync()} + manejo en
\codigo{Proxy.procesarRegistro()}. Membresía persistida: \codigo{RegistroMembresia} en
\codigo{data/MEMBRESIA.txt}. Failover entre nodos del clúster: \codigo{Proxy.rutear()}.

\subsubsection{2.2 Ordenamiento de eventos (Unidad 5) — Relojes de Lamport}
\nivel{Excelente}

\Que Reloj de Lamport thread-safe propagado en todos los mensajes (cliente $\to$ proxy $\to$
servidores, y en Bully/Mutex). Al menos una función muestra \emph{orden causal}: la
conversación de chat se ordena por marca de Lamport. Cada evento relevante queda en el log
con su marca lógica.

\Por Sin reloj global no hay orden temporal absoluto; Lamport da un orden causal
(\emph{happened-before}) suficiente para ordenar la conversación y auditar la coordinación.

\Don \codigo{RelojLamport} (regla \codigo{max(local,recibido)+1}). Estampado en
\codigo{Proxy.rutear()} y \codigo{svMensajeria.enviarMensaje()}; orden causal en
\codigo{svMensajeria.verConversacion()} (campo \codigo{Mensaje.lamportClock}).
Logs: líneas \codigo{[LAMPORT] t=}, \codigo{[BULLY] t=}, \codigo{[MUTEX] t=}.

\begin{lstlisting}[caption={Regla de actualización de Lamport en recepción (com.steam.common.RelojLamport).}]
public long update(long received) {              // evento de recepcion
    long actual, actualizado;
    do {
        actual      = reloj.get();
        actualizado = Math.max(actual, received) + 1;
    } while (!reloj.compareAndSet(actual, actualizado));   // atomico, sin locks
    return actualizado;
}
\end{lstlisting}

\subsubsection{2.3 Coordinación distribuida (Unidad 6) — elegir al menos uno}
\nivel{Excelente — Bully (elección)}

\Que Se implementó el \emph{algoritmo abusón (Bully)} para elegir coordinador entre los
nodos de juegos, que se dispara cuando se detecta la caída del coordinador. Además, ese
coordinador administra la \emph{exclusión mutua} del recurso \codigo{stock} (mutex con
coordinador elegido dinámicamente).

\Por La rúbrica pide al menos un algoritmo de coordinación; Bully cubre la elección de
coordinador ``que se dispare solo cuando se detecta que el coordinador cayó''. El mutex sobre
el stock evita la sobreventa del último ejemplar entre nodos.

\Don Elección: \codigo{GestorBully.iniciarEleccion()/convertirseEnCoordinador()} y su daemon
de heartbeat. Exclusión mutua: \codigo{GestorMutexCentralizado.requestLock()/releaseLock()}
usada en \codigo{svJuegos.comprarJuego()}.
\textit{Nota de defensa:} es exclusión mutua \emph{coordinador-céntrica} (no Ricart--Agrawala);
el requisito se cumple por la elección Bully.

\subsubsection{2.4 Tolerancia a fallos}
\nivel{Excelente}

\Que Detección por \emph{heartbeats y timeouts} en tres niveles y recuperación automática
sin caída total: el Proxy hace health-check (10 s) y conmuta al espejo; Bully envía
heartbeat al coordinador (5 s) y re-elige si cae; el Watchdog reinicia procesos muertos
(15 s, 3 fallos); y un snapshot periódico respalda \codigo{Main}$\to$\codigo{Copy}.

\Por Cubre tanto Crash (nodo caído) como Omisión (mensaje perdido: timeout + reintento), y
garantiza que el servicio sigue tras una caída (re-elección/failover/reinicio).

\Don \codigo{Proxy.iniciarHealthCheck()}; daemon de heartbeat en \codigo{GestorBully};
\codigo{WatchdogServidor}; respaldo en \codigo{GestorSnapshot}.

\subsection{§3 Prueba de Tráfico Real (Carga)}

\subsubsection{3.1 Generador de carga}
\nivel{Excelente}

\Que Cliente Java que lanza 50 hilos concurrentes durante 60 s sostenidos; cada hilo simula
un ciclo de usuario (login $\to$ listar $\to$ comprar $\to$ confirmar $\to$ mensaje $\to$
logout), ejercitando las dos funciones y, en particular, el recurso protegido por el mutex.

\Por La rúbrica exige $\ge$50 clientes/hilos por $\ge$60 s y que la carga golpee el recurso
en exclusión mutua (el camino de compra pide el lock en cada intento).

\Don \codigo{GeneradorCarga} (parámetros por defecto \codigo{50} y \codigo{60}).

\subsubsection{3.2 Métricas recolectadas}
\nivel{Excelente}

\Que Se miden throughput (req/s), latencia promedio y \emph{p95}, cantidad de mensajes del
algoritmo de coordinación, y tasa de pérdida. La métrica de error se corrigió para
distinguir \emph{pérdida real de transporte} (timeout / sin respuesta) de los rechazos de
negocio esperados, de modo que el pico de la falla inducida sea nítido.

\Por La §3.2 pide exactamente esas métricas, incluida la ``cantidad de mensajes que genera
el algoritmo de coordinación''.

\Don Cálculo y reporte en \codigo{GeneradorCarga.imprimirReporteFinal()}. Conteo de mensajes
de coordinación: contadores \codigo{AtomicLong} en \codigo{GestorBully} y
\codigo{GestorMutexCentralizado}, expuestos por la operación \codigo{VER\_METRICAS\_COORD}
de \codigo{svJuegos} y sumados por \codigo{GeneradorCarga.recolectarCoord()}.

\begin{center}\small
\textbf{Última corrida (50 hilos / 60 s, falla del coordinador @30 s):}\\[2pt]
\begin{tabular}{l r}
\hline
Throughput & 412 req/s \\
Latencia p95 & 219 ms \\
Peticiones OK & 16.045 \\
Pérdida (transporte) & 0\% \\
Mensajes de coordinación & 1.964 (Bully = 10, Mutex = 1.954) \\
Tiempo de recuperación tras la falla & $\approx$ 7 s \\
\hline
\end{tabular}
\end{center}

\subsubsection{3.3 Falla inducida durante la prueba}
\nivel{Excelente}

\Que En plena carga se derriba al coordinador a propósito y se mide cuánto tarda el sistema
en reorganizarse. La herramienta localiza al coordinador, le envía un apagado, y mide los
milisegundos hasta que el nodo sobreviviente se proclama nuevo coordinador.

\Por La §3.3 exige derribar al coordinador en vivo y medir recuperación e impacto en
latencia/errores.

\Don \codigo{FallaInducida} (consulta \codigo{QUIEN\_ES\_COORDINADOR}, envía
\codigo{SHUTDOWN\_GRACEFUL} y mide la re-elección). El otro nodo re-elige vía
\codigo{GestorBully}.

\subsubsection{3.4 Evidencia entregada}
\nivel{Logs listos; gráfico en informe}

\Que Se conservan los logs de la corrida (marcas lógicas, eventos de elección y detección de
la falla) y los datos para la tabla de métricas.

\Por La §3.4 exige una tabla y al menos un gráfico, más los logs adjuntos.

\Don Logs: \codigo{logs/carga\_*.log}, \codigo{logs/orq\_falla.out},
\codigo{logs/svJuegos-*\_0.log}, \codigo{data/MEMBRESIA.txt}. Orquestador repetible:
\codigo{scripts/13\_run\_prueba\_trafico.ps1}.
\pend{Gráfico throughput/latencia vs tiempo: informe}

\subsection{§4 Pauta de Evaluación}

\subsubsection{4.1 Teoría — 4.2 Modelado — 4.3 Análisis y Carga (Informe, 35\%)}
\pend{Pendiente: informe}

\Que La sustancia técnica existe: ausencia de reloj global (Lamport), topología multinodo,
modelo de seguridad y de fallos, y la lectura de la prueba de tráfico (Sección §3 de este
documento). Falta redactarlo con el formato pregrado y diagramar el UML.

\Por Es el 35\% del informe; debe interpretar las métricas (throughput, p95, falla inducida),
no solo mostrarlas.

\Don Base para 4.1/4.3 en las secciones §2 y §3 de arriba. \pend{Redacción + diagramas: equipo}

\subsubsection{4.4 Distribución — 4.5 Coordinación y Ordenamiento — 4.6 Tolerancia (Código, 35\%)}
\nivel{Excelente}

\Que Estos tres criterios de implementación corresponden directamente a §2.1, §2.2/§2.3 y
§2.4 ya descritos: topología real + membresía, Lamport + Bully verificables, y
detección/recuperación por heartbeats sin caída total.

\Por Es el núcleo de código del final y donde se concentra la nota de implementación.

\Don Ver §2.1 (\codigo{RegistradorProxy}, \codigo{RegistroMembresia}), §2.2
(\codigo{RelojLamport}), §2.3 (\codigo{GestorBully}, \codigo{GestorMutexCentralizado}) y
§2.4 (\codigo{Proxy.iniciarHealthCheck}, \codigo{WatchdogServidor}).

\subsubsection{4.7 Exposición y Video — 4.8 Demo de Carga y Falla en Vivo — 4.9 Defensa}
\pend{Pendiente: presentación}

\Que La demo de carga + falla en vivo está \emph{lista para ejecutarse} con un solo comando;
solo falta hacerla durante la presentación y grabar el video.

\Por La §4.8 exige correr la prueba de tráfico en vivo sobre $\ge$3 nodos y derribar un nodo
mostrando la recuperación con métricas.

\Don \codigo{scripts/13\_run\_prueba\_trafico.ps1} (levanta los 7 procesos, corre carga + falla
y recoge logs/métricas). \pend{Grabar video $\le$3 min + demo + defensa: equipo}

\subsection{§5 Entregables}

\begin{center}\small
\begin{tabular}{l l l}
\hline
\textbf{Entregable} & \textbf{Estado} & \textbf{Dónde} \\
\hline
Código fuente completo (.zip, incl. generador) & \textcolor{verde}{Listo} & repositorio + \codigo{sistema-steam.jar} \\
Logs de carga y falla inducida & \textcolor{verde}{Listo} & \codigo{logs/} \\
Informe PDF (formato pregrado + portada) & \textcolor{acento}{Pendiente} & equipo \\
Presentación con video demostrativo & \textcolor{acento}{Pendiente} & equipo \\
\hline
\end{tabular}
\end{center}

% ═════════════════════════════════════════════════════════════════════════════
\section{Apéndice — Puntos a cuidar en la defensa}
% ═════════════════════════════════════════════════════════════════════════════

\begin{itemize}[leftmargin=1.4em]
  \item \textbf{El Proxy es un punto único de \emph{ruteo}}, pero la \emph{coordinación} (Bully,
        Mutex) está distribuida entre los nodos de juegos, no centralizada. Conviene decirlo así.
  \item \textbf{El mutex es coordinador-céntrico} (no Ricart--Agrawala): preséntenlo como
        ``exclusión mutua administrada por el coordinador elegido dinámicamente por Bully''.
        El requisito §2.3 se cumple por la \emph{elección} Bully.
  \item \textbf{Reloj de Lamport (escalar)}: cumple §2.2. Si preguntan por vectoriales, la
        respuesta es que Lamport basta para el orden total del chat que se requiere.
  \item \textbf{Recuperación medida ($\approx$7 s, 0\% de pérdida)}: el sistema sigue
        sirviendo durante la caída del coordinador gracias al failover del Proxy al nodo espejo.
\end{itemize}

\end{document}
